\begin{abstract}
Geostationary imagers see high clouds displaced away from the sub-satellite point. For INSAT-3DR at 74\,\textdegree E, deep convection over the Himalaya and Northeast India appears about 10\,km from its true position, two to three infrared pixels. We present PyGeoParallax-IN, an open-source Python tool that corrects this parallax for INSAT-3DR imagery. It estimates cloud-top height from 10.8\,\textmu m brightness temperature, traces each line of sight to the cloud top on the WGS84 ellipsoid (with terrain if a digital elevation model is given), applies a navigation correction measured from coastlines, and regrids the scene so that uncorrected and corrected datasets share one grid.

We validate the correction against GPM Ku-band radar in 33 coincident monsoon cases (2022--2024) over South, Central and Northeast India and the Himalaya, 14 of them usable. Parallax correction lowers the median INSAT--GPM storm-top offset from 5.5 to 3.1\,km (10 of 12 cases improve; sign test $p=0.039$) and leaves the cross-look offset unchanged. At storm cores, however, a parallax-only correction overshoots by about 5.7\,km: INSAT-3DR L1C positions carry a seasonal navigation error, 7--20\,km to the west in winter and 2.5--5\,km to the south in the monsoon, which partly cancels the parallax. Correcting navigation as well reduces the overshoot significantly and gives a median field offset of 1.8\,km, although with 12 cases this is not significantly better than parallax correction alone. ERA5 profiles and SRTM terrain change the storm results by less than 0.5\,km. The tool processes an India scene in 1.8\,s.
\end{abstract}
